\documentclass[11pt]{article}
\usepackage[margin=1in]{geometry}
\usepackage[T1]{fontenc}
\usepackage{lmodern}
\usepackage{enumitem}
\usepackage{xcolor}
\usepackage{hyperref}
\usepackage{listings}

\hypersetup{colorlinks=true, linkcolor=blue, urlcolor=blue}
\lstset{
  basicstyle=\ttfamily\small,
  columns=fullflexible,
  keepspaces=true,
  frame=single,
  breaklines=true
}

\title{Lab 2: Memory Hierarchy\\Progress Summary}
\author{}
\date{October 21, 2026}

\begin{document}
\maketitle

\section{Lab premise}

Lab 2 asks us to build a software model of a processor memory hierarchy and use it
to study cache design. The long-term goal is to choose a practical two-level
cache organization that provides low Average Memory Access Time (AMAT) while
keeping area and power reasonable.

The lab progresses through four stages:
\begin{enumerate}[leftmargin=*]
  \item Implement a configurable set-associative cache.
  \item Compare random, LRU, bit-PLRU, and other replacement policies.
  \item Build a parameterized L1/L2 cache hierarchy.
  \item Search the design space under a total 5 MB cache-storage budget, using
        benchmark results and Destiny hardware-cost estimates.
\end{enumerate}

The central tradeoff is that larger blocks, greater associativity, and larger
caches can reduce misses, but can also increase access latency, area, leakage
power, dynamic power, and tag-storage overhead.

\section{Repository and execution model}

The project contains a Pin-based instrumentation framework:
\begin{itemize}[leftmargin=*]
  \item \texttt{memory-hierarchy.cpp} instruments every executed memory
        operand. Reads call \texttt{mhLoad}; writes call \texttt{mhStore}.
  \item \texttt{cache.h} declares the cache model and the interface called by
        the Pin instrumentation.
  \item \texttt{cache.cpp} implements the cache model.
  \item \texttt{makefile} and \texttt{makefile.rules} describe how the Pin
        tool is built.
  \item The PDF writeup describes the required experiments and analysis.
\end{itemize}

The runtime flow is:
\begin{verbatim}
benchmark instruction
        -> Pin memory instrumentation
        -> mhLoad() or mhStore()
        -> simulated cache
        -> statistics
        -> memhier.stats at program termination
\end{verbatim}

\section{Work completed}

\subsection{\texttt{cache.h}}

The header now defines a first single-cache abstraction:
\begin{itemize}[leftmargin=*]
  \item \texttt{CacheLine}, containing a valid bit and tag;
  \item \texttt{CacheStats}, containing accesses, hits, and misses;
  \item \texttt{Cache}, parameterized by capacity, block size, and
        associativity;
  \item vectors representing sets and ways;
  \item a per-set victim pointer for temporary round-robin replacement;
  \item \texttt{mhInit()} and \texttt{mhPrintStats()} for Pin integration;
  \item the existing \texttt{mhLoad()} and \texttt{mhStore()} entry points.
\end{itemize}

\subsection{\texttt{cache.cpp}}

The first working cache implementation is complete. It performs the following:
\begin{enumerate}[leftmargin=*]
  \item Validates cache dimensions.
  \item Computes the number of blocks and sets.
  \item Allocates each set with the requested number of ways.
  \item Converts an address to a block address.
  \item Computes the set index and tag:
\begin{verbatim}
block address = address / block size
set index     = block address % number of sets
tag           = block address / number of sets
\end{verbatim}
  \item Searches every way in the selected set for a matching valid tag.
  \item Records hits and misses.
  \item Uses an invalid way before evicting a valid line.
  \item Uses simple round-robin replacement when a set is full.
  \item Splits an access that crosses a block boundary into multiple block
        accesses.
  \item Prints accesses, hits, misses, and miss rate.
\end{enumerate}

For example, an 8-byte access beginning at byte 60 in a cache with 16-byte
blocks touches two simulated cache blocks.

\subsection{\texttt{memory-hierarchy.cpp}}

The existing Pin instrumentation was preserved. The following integration was
added:
\begin{itemize}[leftmargin=*]
  \item After Pin parses command-line arguments, \texttt{mhInit()} is called.
  \item \texttt{-l1c} is converted from kilobytes to bytes.
  \item \texttt{-l1b} supplies the block size.
  \item \texttt{-l1a} supplies the associativity.
  \item At program termination, \texttt{Fini()} calls \texttt{mhPrintStats()}.
\end{itemize}

The current default configuration is equivalent to:
\begin{lstlisting}
capacity:      2 KB
block size:    64 bytes
associativity: 2-way
\end{lstlisting}

\section{Validation performed}

The implementation was compiled independently of Pin with:
\begin{lstlisting}
c++ -std=c++11 -Wall -Wextra -pedantic -I. cache.cpp cache_test.cpp
\end{lstlisting}

A deterministic test exercised:
\begin{itemize}[leftmargin=*]
  \item a repeated access to the same block;
  \item an access within an already cached block;
  \item a store crossing a block boundary;
  \item the \texttt{mhInit()} and \texttt{mhPrintStats()} interface.
\end{itemize}

The resulting statistics were:
\begin{lstlisting}
accesses: 4
hits: 1
misses: 3
miss_rate: 0.75
\end{lstlisting}

This confirms that the cache can be configured through the public interface,
that repeated block accesses hit, and that a boundary-crossing access is
counted once for each block touched.

\section{Current status and limitations}

The project currently has a correct first single-cache milestone, but it is not
yet the completed lab. The following work remains:
\begin{itemize}[leftmargin=*]
  \item Run the actual Pin tool on a small program and verify
        \texttt{memhier.stats}.
  \item Replace the temporary round-robin policy with random, LRU, and bit-PLRU.
  \item Add load/store policy decisions such as write-back or write-through and
        write-allocate behavior.
  \item Add L2 and DRAM behavior.
  \item Add latency and AMAT accounting.
  \item Connect Destiny area, leakage-power, dynamic-power, and latency data.
  \item Build benchmark automation, design-space search, Pareto-frontier plots,
        and sensitivity studies.
\end{itemize}

The actual Pin build/run has not yet been performed in this work session because
the Pin kit location and executable test workload have not been configured.
The cache implementation has instead been validated with a standalone C++
test.

\section{Recommended next step}

The next task is to verify the instrumented Pin path with a very small test
program. After that, use the existing L1 knobs for experiments. Only once the
single-cache hit/miss behavior is trusted should replacement policies, L2, AMAT,
and design-space exploration be added.

\end{document}
